\chapter*{Введение}
\addcontentsline{toc}{chapter}{Введение}

\section*{Актуальность темы}

Современная биоинформатика работает с данными огромного масштаба. Геномы, транскриптомы и белковые последовательности становятся всё доступнее, но способы их восприятия остаются преимущественно визуальными и числовыми. Для специалиста последовательность из миллионов символов является материалом для вычислительного анализа, требующим сжатия, сравнения и статистической обработки.

Музыкальная сонификация предлагает альтернативный способ взаимодействия с биологическими данными. Человек умеет воспринимать сложные временные структуры не только зрением, но и слухом. Музыка естественным образом организует высоту, длительность, повтор, плотность и динамику. Поэтому возникает вопрос: можно ли использовать музыку как дополнительный язык восприятия биологических последовательностей?

Актуальность работы определяется несколькими факторами. Во-первых, объём биологических данных быстро растёт, требуя новых способов компактного представления. Во-вторых, сонификация является перспективным методом восприятия научных данных, раскрывающим временные изменения иначе, чем визуализация. В-третьих, генеративные модели позволяют уйти от прямолинейного отображения символов в ноты, создавая музыкально связные композиции под управлением биологических признаков.

\section*{Степень разработанности темы}

Задача сонификации биологических данных исследуется с 1980-х годов. Ранние работы использовали прямое отображение нуклеотидов в ноты~\cite{ohno1986}. Современные подходы включают алгоритмическую композицию~\cite{temple2017}, параметрическую сонификацию~\cite{grond2011} и генеративные модели.

В области генерации символической музыки активно развиваются методы на основе рекуррентных сетей~\cite{eck2002}, вариационных автоэнкодеров~\cite{roberts2018}, генеративно-состязательных сетей~\cite{dong2018} и трансформеров~\cite{huang2018}. Обзорные работы по deep music generation подчёркивают, что качество результата зависит не только от мощности модели, но и от выбранного музыкального представления, метрик и устойчивости обучения~\cite{briot2017,ji2020}. Датасеты POP909~\cite{wang2020} и MAESTRO~\cite{hawthorne2019} стали стандартом для обучения моделей.

В предварительных исследованиях перед выполнением работы также рассматривались вопросы выбора музыкальной нотации и оценки результата: MusicLang как высокоуровневый язык описания музыки, MIDI как универсальный событийный формат, а также идея модели-оценщика для музыкальных фрагментов. Эти исследования привели к выводу, что для финальной системы важнее не прямое кодирование символов, а промежуточное структурированное представление, пригодное для обучения и последующей количественной оценки.

Однако большинство существующих решений либо используют слишком простое отображение символов, либо не учитывают музыкальную структуру. Задача создания системы, которая генерирует музыкально связные композиции под управлением биологических признаков, остаётся актуальной.

\section*{Проблема}

Существующие подходы к биосонификации имеют ограничения:
\begin{itemize}
\item Прямое отображение символов в ноты даёт музыкально бедный результат
\item Алгоритмические методы требуют ручной настройки правил
\item Генеративные модели без структурного разделения создают хаотичные последовательности
\item Отсутствуют воспроизводимые эксперименты с оценкой качества
\end{itemize}

\section*{Объект и предмет исследования}

\textbf{Объект исследования} — биологические последовательности и способы их преобразования в воспринимаемую человеком форму.

\textbf{Предмет исследования} — метод нейросетевой генерации символической музыки, условно управляемой признаками биологических последовательностей.

\section*{Цель работы}

Разработать и экспериментально проверить обучаемую нейросетевую модель для синтеза и оценки мелодий по биологическим данным, используя биологические признаки как управляющий сигнал для нейросетевой генерации.

\section*{Уточнение термина ``самообучающаяся модель''}

В утверждённой теме работы используется термин ``самообучающаяся модель''. На начальном этапе проектирования рассматривалась генеративно-состязательная схема, в которой одна модель генерирует музыкальный материал, а другая модель-оценщик формирует обучающий сигнал для улучшения генератора. Такая постановка соответствует идее итеративного улучшения за счёт внутренней оценки результата.

В ходе анализа предметной области и вычислительных ограничений итоговая архитектура была изменена. Для задачи генерации символической музыки по биологическим данным была выбрана более устойчивая и воспроизводимая схема: условный autoregressive Transformer, обучаемый на парах ``биологический фрагмент — музыкальный сегмент''. Поэтому в дальнейшем термин ``самообучающаяся'' трактуется не как автономное онлайн-обучение в процессе работы веб-интерфейса и не как финальная GAN-архитектура, а как обучаемая на данных нейросетевая модель, параметры которой автоматически оптимизируются по функции потерь на обучающей выборке. Оценка качества вынесена в экспериментальный контур: модель сравнивается с baseline и предыдущей версией по заранее заданным метрикам.

\section*{Задачи работы}

Для достижения цели были поставлены следующие задачи:
\begin{enumerate}
\item Изучить предметную область биосонификации и генерации символической музыки
\item Проанализировать существующие решения и выявить их недостатки
\item Подготовить биологические данные (FASTA) и музыкальный корпус (MIDI)
\item Реализовать извлечение биологических признаков из DNA, RNA и protein последовательностей
\item Разработать иерархическую архитектуру генерации: Bio→Harmony→Melody
\item Реализовать pairing биологических и музыкальных сегментов
\item Обучить модели на доступном оборудовании (NVIDIA RTX 2060 6GB)
\item Провести экспериментальную оценку и сравнение с baseline
\item Подготовить воспроизводимые результаты и документацию
\end{enumerate}

\section*{Методы исследования}

В работе использовались следующие методы:
\begin{itemize}
\item Анализ научной литературы и существующих решений
\item Проектирование программной архитектуры
\item Реализация программного комплекса на Python
\item Машинное обучение с использованием PyTorch
\item Экспериментальная оценка на тестовых данных
\item Сравнительный анализ с baseline
\end{itemize}

\section*{Научная новизна}

Научная новизна работы состоит в предложении иерархической архитектуры, где биологические признаки используются как управляющий сигнал для структурированной музыкальной генерации. В отличие от прямого отображения символов, система разделяет задачу на два этапа: генерацию гармонии и генерацию мелодии поверх гармонии. Это обеспечивает музыкальную связность при сохранении управления от биологических данных.

\section*{Практическая значимость}

Практическая значимость работы состоит в создании рабочего программного прототипа BioSonification. Система включает:
\begin{itemize}
\item Обученные модели с воспроизводимыми checkpoint
\item Веб-интерфейс для генерации музыки из FASTA
\item Документацию и примеры использования
\item Отчёты по данным и метрикам качества
\end{itemize}

Система может использоваться в образовательных целях, для научной коммуникации и демонстрации биоинформатических данных.

\section*{Структура работы}

Работа состоит из введения, трёх глав, заключения, списка литературы и приложений.

В \textbf{первой главе} анализируется предметная область, существующие подходы к биосонификации и генерации музыки, выявляются недостатки аналогов и формулируется постановка задачи.

Во \textbf{второй главе} описывается проектирование и реализация системы: техническое задание, архитектура, биологический энкодер, музыкальное представление, pairing, нейросетевая модель и процесс обучения.

В \textbf{третьей главе} представлена экспериментальная оценка: методика проверки, метрики, результаты обучения, сравнение с baseline и анализ ограничений.

В \textbf{заключении} подводятся итоги работы и намечаются направления дальнейшего развития.
